\begin{helper}
On the actual split-outside-blocker witness surface, every adjacent cut
\(k<n\) has a genuine two-piece decomposition of the complete source and
target stages.  Fix the activation--priority measures \(\nu,\nu'\), the local embedding
\(\phi\), the private blocker map \(\beta\), the injective enumeration
\(z_0,\ldots,z_{n-1}\), the candidate sets
\[
B_k=\{x\in X:xz_k\in E(G)\},
\]
the paired atom families \(C_\omega,C'_\omega\), the common nonnegative
weights \(w(\omega)\), and the actual one-changing-coordinate boundary events
\(L_k(\omega)\) and \(R_k(\omega)\).
Assume that these boundary events have already been registered as genuine
subevents of the corresponding complete stages, after the actual stage
definitions are expanded, and that the boundary-cell intersections
\[
L_k(\omega)\cap C_\omega,\qquad R_k(\omega)\cap C'_\omega
\]
are measurable in their source and target activation--priority product
spaces.  These are the side conditions supplied by the measured-stage
refinement layer; this two-piece node is only the sidewise partition and
normalization step.

For each \(k<n\) and \(\omega\), let \(S_k(\omega)\) be the source complete
stage in which a candidate \(x\in X\) is active, beats every active embedded
neighbor in \(X\cup\{r\}\), and is not killed by any retained common blocker
\(z_i\) with \(i\ge k\) and \(x\in B_i\).  Let \(T_{k+1}(\omega)\) be the
target complete stage in which \(\phi(x)\) is active, beats the embedded
neighbors, and is not killed by any private blocker already installed at an
index \(i<k+1\).

Then there are source and target residual unchanged events
\(U^-_k(\omega)\), \(U^+_k(\omega)\), boundary events
\(D^-_k(\omega)=L_k(\omega)\), \(D^+_k(\omega)=R_k(\omega)\), and normalized
real representatives
\[
H^-_k,\ U^-_k,\ D^-_k,\qquad H^+_{k+1},\ U^+_k,\ D^+_k
\]
such that \(U^-_k(\omega)\) is the complement of \(D^-_k(\omega)\) inside
\(S_k(\omega)\), and \(U^+_k(\omega)\) is the complement of
\(D^+_k(\omega)\) inside \(T_{k+1}(\omega)\).  Thus, after intersecting with
the prescribed cells,
\[
S_k(\omega)=(U^-_k(\omega)\cup D^-_k(\omega)),\qquad
T_{k+1}(\omega)=(U^+_k(\omega)\cup D^+_k(\omega)),
\]
the two pieces on each side are disjoint modulo the cell, the measure of the
complete stage is the sum of the two piece measures, all representatives are
nonnegative, and if \(w(\omega)=0\) then all six representatives are zero.
Moreover these real representatives are genuine normalized representatives of
the three subevents on each side:
\[
\operatorname{ofReal}\!\bigl(w(\omega)H^-_k(\omega)\bigr)
 =\nu(S_k(\omega)\cap C_\omega),\quad
\operatorname{ofReal}\!\bigl(w(\omega)U^-_k(\omega)\bigr)
 =\nu(U^-_k(\omega)\cap C_\omega),
\]
\[
\operatorname{ofReal}\!\bigl(w(\omega)D^-_k(\omega)\bigr)
 =\nu(D^-_k(\omega)\cap C_\omega),
\]
and
\[
\operatorname{ofReal}\!\bigl(w(\omega)H^+_{k+1}(\omega)\bigr)
 =\nu'(T_{k+1}(\omega)\cap C'_\omega),\quad
\operatorname{ofReal}\!\bigl(w(\omega)U^+_k(\omega)\bigr)
 =\nu'(U^+_k(\omega)\cap C'_\omega),
\]
\[
\operatorname{ofReal}\!\bigl(w(\omega)D^+_k(\omega)\bigr)
 =\nu'(D^+_k(\omega)\cap C'_\omega).
\]
Consequently,
\[
H^-_k(\omega)=U^-_k(\omega)+D^-_k(\omega),\qquad
H^+_{k+1}(\omega)=U^+_k(\omega)+D^+_k(\omega),
\]
with the displayed boundary pieces equal to the actual adjacent boundary
events \(L_k(\omega)\) and \(R_k(\omega)\).
\end{helper}
\begin{proof}
Fix \(k<n\) and \(\omega\).  The actual stage surface is the one provided by
\(\noderef{SamplingSplitOutsideBlockerTrueHybridStageSurface}\): at cut \(j\)
the target has private tests for the indices \(i<j\), while the source
retains common blocker tests for the indices \(j\le i<n\).  Thus the source
complete stage at cut \(k\) consists of the embedded-neighbor survival core
together with all retained tests \(i\ge k\), and the target complete stage at
cut \(k+1\) consists of the transported embedded-neighbor core together with
all private tests \(i<k+1\).

Separate the current boundary event from those two complete stages.  On the
source side, the boundary event \(D^-_k(\omega)\) is exactly
\(L_k(\omega)\): it records a witnessing candidate \(x\in B_k\) which
satisfies the embedded-neighbor core, all retained common-blocker tests
except the current index, and also satisfies the current common-blocker
survival test at \(z_k\).  The unchanged source event is not the larger
omitted-current context, since that larger context contains the boundary
event.  Instead set
\[
U^-_k(\omega)=S_k(\omega)\setminus D^-_k(\omega).
\]
Then \(S_k(\omega)=U^-_k(\omega)\cup D^-_k(\omega)\) and the two pieces are
disjoint, both before and after intersection with \(C_\omega\).  This is the
genuine partition of the cut-\(k\) complete source stage into the current
boundary part and its residual nonboundary part.

The target side is identical with the roles transported through the private
blocker map.  At cut \(k+1\) the current private boundary event is
\(D^+_k(\omega)=R_k(\omega)\), identified using the supplied flat
representative of the dependent private vertex.  Define
\[
U^+_k(\omega)=T_{k+1}(\omega)\setminus D^+_k(\omega).
\]
Then \(T_{k+1}(\omega)=U^+_k(\omega)\cup D^+_k(\omega)\), and the two target
pieces are disjoint on \(C'_\omega\).  Thus both sides use the same kind of
residual-inside-complete partition; neither side assumes disjointness for a
nested omitted-current event.

The finite-additivity identities now follow from the displayed set
identities and the stated measured-stage side conditions.  Indeed, after
intersecting with \(C_\omega\), the source complete event is exactly the
union of \(U^-_k(\omega)\cap C_\omega\) and
\(D^-_k(\omega)\cap C_\omega\).  The reverse containment uses the explicit
assumption that \(L_k(\omega)=D^-_k(\omega)\) is a subevent of
\(S_k(\omega)\).  The two pieces are disjoint because
\(U^-_k(\omega)=S_k(\omega)\setminus D^-_k(\omega)\), and the second piece
\(D^-_k(\omega)\cap C_\omega\) is measurable by the hypothesis above.
Finite additivity for a disjoint union with measurable second piece gives
\[
\nu(S_k(\omega)\cap C_\omega)
=\nu(U^-_k(\omega)\cap C_\omega)+\nu(D^-_k(\omega)\cap C_\omega)
\]
on the source side.  The same argument applies to the target side, using the
assumed containment \(R_k(\omega)\subseteq T_{k+1}(\omega)\) and the
measurability of \(R_k(\omega)\cap C'_\omega\), and gives
\[
\nu'(T_{k+1}(\omega)\cap C'_\omega)
=\nu'(U^+_k(\omega)\cap C'_\omega)+\nu'(D^+_k(\omega)\cap C'_\omega).
\]
For each of the six subevents choose the normalized real representative by
\(\noderef{SamplingFiniteMassSubcellNormalization}\).  The hypotheses of
that normalization lemma are checked separately for each subevent.  The
common weight is nonnegative by the theorem hypothesis.  The cell mass
identity is
\(\operatorname{ofReal}(w(\omega))=\nu(C_\omega)\) on the three source
subevents and
\(\operatorname{ofReal}(w(\omega))=\nu'(C'_\omega)\) on the three target
subevents.  Finally each subevent after cell intersection is a subset of
the relevant cell: for example
\(S_k(\omega)\cap C_\omega\subseteq C_\omega\),
\(U^-_k(\omega)\cap C_\omega\subseteq C_\omega\), and
\(D^-_k(\omega)\cap C_\omega\subseteq C_\omega\), and similarly on the
target side.  Therefore the six applications of the normalization lemma
produce nonnegative representatives satisfying the six displayed
\(\operatorname{ofReal}(w\cdot\hbox{mass})\) equations, and in the branch
\(w(\omega)=0\) the same lemma gives the stated convention that all six
representatives are \(0\).

Finally apply the one-sided algebra lemma
\(\noderef{SamplingSplitOutsideBlockerMixedStageCompleteDecomposition}\)
first to the source measure, source cell, cut-\(k\) complete stage,
source residual unchanged event, and source boundary event, and then to the
target measure, target cell, cut-\((k+1)\) complete stage, target residual
unchanged event, and target boundary event.  The hypotheses just established
are exactly its cover, disjointness, finite-additivity, nonnegativity,
zero-weight, and three mass-representation inputs.  It yields
\[
H^-_k(\omega)=U^-_k(\omega)+D^-_k(\omega),\qquad
H^+_{k+1}(\omega)=U^+_k(\omega)+D^+_k(\omega).
\]
Since \(k\) and \(\omega\) were arbitrary, the families can be assembled for
all adjacent cuts and all atoms on the single witness surface.
\end{proof}
